%%%%%%%%%%%%%%%%
%%% deut 32 %%%
%%%%%%%%%%%%%%%%

\Chapter{32}

\Verse{1}
{
	\hProto{הַאֲזִינוּ הַשָּׁמַיִם וַאֲדַבֵּרָה וְתִשְׁמַע הָאָרֶץ אִמְרֵי פִי׃}
}
{
	\eDtrB{
		LISTEN 32 Listen, skies, so I may speak and let the earth
		hear what my mouth says.
	}
}
{
}

\Verse{2}
{
	\hProto{יַעֲרֹף כַּמָּטָר לִקְחִי תִּזַּל כַּטַּל אִמְרָתִי כִּשְׂעִירִם עֲלֵי דֶשֶׁא וְכִרְבִיבִים עֲלֵי עֵשֶׂב׃}
}
{
	\eDtrB{
		Let my teaching come down like showers; let my saying
		emerge like dew, like raindrops on plants and like
		rainfalls on herbs.
	}
}
{
%	\footnote{
%		plants. This word, too, occurs among the first words of the Torah and
%		does not occur again until its appearance here. In Genesis, God causes
%		the earth to produce plants by His word (Gen 1:11–12). Now Moses turns
%		this image and sings of his words dropping like rain on the plants.
%	}
}

\Verse{3}
{
	\hProto{כִּי שֵׁם יְהֹוָה אֶקְרָא הָבוּ גֹדֶל לֵאלֹהֵינוּ׃}
}
{
	\eDtrB{When I call YHWH’s name, avow our {\God}’s greatness.}
}
{
}

\Verse{4}
{
	\hProto{הַצּוּר תָּמִים פּעֳלוֹ כִּי כל דְּרָכָיו מִשְׁפָּט אֵל אֱמוּנָה וְאֵין עָוֶל צַדִּיק וְיָשָׁר הוּא׃}
}
{
	\eDtrB{
		The Rock: His work is unblemished, for all His ways are
		judgment. A {\God} of trust, and without injustice, He’s
		virtuous and right.
	}
}
{
%	\footnote{
%		The Rock. The first letter of the word is oversized in the Torah,
%		possibly to distinguish it from other uses in the song. The song plays
%		on this word, using it seven more times with several different meanings.
%	}
}

\Verse{5}
{
	\hProto{שִׁחֵת לוֹ לֹא בָּנָיו מוּמָם דּוֹר עִקֵּשׁ וּפְתַלְתֹּל׃}
}
{
	\eDtrB{
		It corrupted at Him—not His children, their flaw— a
		crooked and twisted generation.
	}
}
{
%	\footnote{
%		It corrupted at Him—not His children, their flaw. The noun, Hebrew mûm,
%		means an injury of the sort that makes a man unacceptable for the
%		priesthood, or an animal unacceptable for sacrifice. Here I understand
%		it to be some such damage in humans’ nature, as God and Moses have just
%		said (31:21,27,29). It is not humans, “His children,” as such, but
%		rather this inherent shortcoming in their nature that brings about
%		corruption in the eyes of the deity.
%	}
}

\Verse{6}
{
	\hProto{הַ לְיְהֹוָה תִּגְמְלוּ זֹאת עַם נָבָל וְלֹא חָכָם הֲלוֹא הוּא אָבִיךָ קָּנֶךָ הוּא עָשְׂךָ וַיְכֹנְנֶךָ׃}
}
{
	\eDtrB{
		Is it to YHWH that you repay like this?! Foolish people
		and unwise! Isn’t He your father, who created you, He who
		made you and reared you?
	}
}
{
%	\footnote{
%		Is it to YHWH … ?! The first letter (the Hebrew interrogative h, which
%		makes the line a question) is oversized in many manuscripts. This seems
%		to underscore the question, to make it incredible: “You’d repay God this
%		way?!”
%	}
}

\Verse{7}
{
	\hProto{זְכֹר יְמוֹת עוֹלָם בִּינוּ שְׁנוֹת דֹּר וָדֹר שְׁאַל אָבִיךָ וְיַגֵּדְךָ זְקֵנֶיךָ וְיֹאמְרוּ לָךְ׃}
}
{
	\eDtrB{
		Remember the days of old. Grasp the years through
		generations. Ask your father, and he’ll tell you, your
		elders, and they’ll say to you:
	}
}
{
}

\Verse{8}
{
	\hProto{בְּהַנְחֵל עֶלְיוֹן גּוֹיִם בְּהַפְרִידוֹ בְּנֵי אָדָם יַצֵּב גְּבֻלֹת עַמִּים לְמִסְפַּר בְּנֵי יִשְׂרָאֵל׃}
}
{
	\eDtrB{
		When the Highest gave nations legacies, when He dispersed
		humankind, He set the peoples’ borders to the number of
		the children of Israel.
	}
}
{
%	\footnote{
%		dispersed humankind. The term “dispersed” refers to the separation of
%		humankind into nations in Genesis (10:5,18,32; 25:23). Footnote 32:8. to
%		the number of the children of Israel. In light of the other occurrences
%		of the phrase “to the number” in the Tanak (Josh 4:5,8; Judg 21:23), it
%		means here “equal to the number of the children of Israel.” This makes
%		little sense: that the number of nations is equal to the number of the
%		children of Israel. Traditional commentators have taken this to mean
%		that there are seventy nations (as in Genesis 10), matching the seventy
%		souls of Jacob’s (Israel’s) family who go to Egypt (Exod 1:5), or that
%		there are twelve Canaanite peoples, matching Israel’s twelve sons.
%		Contemporary scholars have taken account of a Qumran (Dead Sea) scroll
%		(4QDeut\textasciicircum{}(j)) that says instead: “to the number of the
%		children of the gods” (or “children of God”—b[image]nê
%		’[image]l[image]hîm). The phrase “children [or: sons] of the gods”
%		normally means: the gods. Scholars have hypothesized that this was the
%		original reading of this verse and that later scribes or editors were
%		distressed by this meaning so they changed “children of the gods” to
%		“children of Israel.” If so, then the original meaning was that God
%		allotted a nation to each of the gods—and He took Israel as His own
%		people. Thus the next verse reads: “YHWH’s portion is Israel.” One may
%		take this to mean that the author of this song believed that there were
%		in fact other gods but that YHWH was supreme. (This is known as
%		henotheism, a step between polytheism and monotheism.) Or one may
%		understand it to mean that the author’s view was that YHWH allotted the
%		peoples their gods in whom to believe, but that these gods did not
%		really exist. To complicate this picture further, the Greek text reads
%		“to the number of the angels of God.” And there are cases in which the
%		b[image]nê ’[image]l[image]hîm are thought to be the angels (Gen 6:2;
%		Job 1:6). It is likely that the Greek translator had the “children of
%		the gods” reading and changed it to “angels of God” to establish that
%		these were in fact the angels and not the pagan gods. On this
%		understanding of the b[image]nê ’[image]l[image]hîm, God allots a nation
%		to each of the angels but retains Israel for Himself. I believe that
%		Psalm 82 is particularly relevant to this matter. It pictures the
%		following: “God is standing in the divine council. He judges among the
%		gods” (82:1). In this extraordinary scene, God criticizes the gods for
%		failing to judge justly, and He takes away their immortality, condemning
%		them to die like humans: “I had said you are gods, children of the
%		Highest, all of you; but you shall die like a human!” (82:6). Whether
%		meant literally or figuratively, this psalm tells a myth of the death of
%		the gods. And note that it refers to God as “the Highest”
%		(‘ely[image]n), as in our verse here in the Song of Moses (and it refers
%		to the gods as the children of the Highest); and it ends with God giving
%		legacies to the nations (82:8)—which is also what our verse here says:
%		“When the Highest gave nations legacies....” It is likely, therefore,
%		that the passage in the Song of Moses reflects the idea of Psalm 82:
%		that there once were lesser gods along with YHWH, and that each was the
%		god of a people, but they were inadequate, and they no longer exist.
%		This may have been believed literally, or it may have been conceived as
%		an answer to those who asked, “What happened to the gods we used to
%		worship?” Or it may have been conceived as a metaphor for the
%		replacement of the belief in many gods by the belief in one.
%	}
}

\Verse{9}
{
	\hProto{כִּי חֵלֶק יְהֹוָה עַמּוֹ יַעֲקֹב חֶבֶל נַחֲלָתוֹ׃}
}
{
	\eDtrB{
		For YHWH’s portion is His people. Jacob is the share of
		His legacy.
	}
}
{
}

\Verse{10}
{
	\hProto{יִמְצָאֵהוּ בְּאֶרֶץ מִדְבָּר וּבְתֹהוּ יְלֵל יְשִׁמֹן יְסֹבְבֶנְהוּ יְבוֹנְנֵהוּ יִצְּרֶנְהוּ כְּאִישׁוֹן עֵינוֹ׃}
}
{
	\eDtrB{
		He found it in a wilderness land and in a formless place,
		a howling desert. He surrounded it. He attended to it. He
		guarded it, like the pupil of his eye.
	}
}
{
%	\footnote{
%		a formless place. Hebrew tôhû. The word has not occurred since its
%		mention at the very beginning of the Torah: “the earth had been formless
%		and shapeless” (Gen 1:2). Now Israel’s environment in the wilderness is
%		pictured comparably to the condition of the universe prior to the acts
%		of creation: a condition of chaos.
%	}
}

\Verse{11}
{
	\hProto{כְּנֶשֶׁר יָעִיר קִנּוֹ עַל גּוֹזָלָיו יְרַחֵף יִפְרֹשׂ כְּנָפָיו יִקָּחֵהוּ יִשָּׂאֵהוּ עַל אֶבְרָתוֹ׃}
}
{
	\eDtrB{
		As an eagle stirs its nest, hovers over its young, spreads
		its wings, takes it, lifts it on its pinion,
	}
}
{
%	\footnote{
%		hovers. This word, too, occurred in Gen 1:2 (see the preceding comment)
%		and never again in the Torah until its appearance here. YHWH’s
%		protection of Israel in its infancy in the wilderness (like an eagle
%		hovering over its young) is pictured comparably to the divine spirit
%		hovering over the still-shapeless waters in which the earth and skies
%		will be formed.
%	}
}

\Verse{12}
{
	\hProto{יְהֹוָה בָּדָד יַנְחֶנּוּ וְאֵין עִמּוֹ אֵל נֵכָר׃}
}
{
	\eDtrB{YHWH, alone, led it, and no foreign god with Him.}
}
{
}

\Verse{13}
{
	\hProto{יַרְכִּבֵהוּ עַל [בָּמֳתֵי] (במותי) אָרֶץ וַיֹּאכַל תְּנוּבֹת שָׂדָי וַיֵּנִקֵהוּ דְבַשׁ מִסֶּלַע וְשֶׁמֶן מֵחַלְמִישׁ צוּר׃}
}
{
	\eDtrB{
		He had it ride over earth’s high places and fed it the
		field’s bounties and had it suck honey from a rock and oil
		from a flint rock
	}
}
{
}

\Verse{14}
{
	\hProto{חֶמְאַת בָּקָר וַחֲלֵב צֹאן עִם חֵלֶב כָּרִים וְאֵילִים בְּנֵי בָשָׁן וְעַתּוּדִים עִם חֵלֶב כִּלְיוֹת חִטָּה וְדַם עֵנָב תִּשְׁתֶּה חָמֶר׃}
}
{
	\eDtrB{
		and curds of cattle and milk of the flock and fat of lambs
		and Bashan rams and he-goats with fat of innards of wheat
		—and from grape’s blood you drank wine.
	}
}
{
}

\Verse{15}
{
	\hProto{וַיִּשְׁמַן יְשֻׁרוּן וַיִּבְעָט שָׁמַנְתָּ עָבִיתָ כָּשִׂיתָ וַיִּטֹּשׁ אֱלוֹהַּ עָשָׂהוּ וַיְנַבֵּל צוּר יְשֻׁעָתוֹ׃}
}
{
	\eDtrB{
		And Jeshurun got fat and kicked —you got fat, you got
		wide, you got stuffed!— and it left {\God} who made it and
		took its saving rock for granted.
	}
}
{
%	\footnote{
%		Jeshurun. Israel. Footnote 32:15. got fat. This is a pun. The Hebrew
%		consonants y[image]mn are the same as in the word “desert”
%		(y[image][image]im[image]n) in v. 10. But the wordplay is that the same
%		words consonantally have nearly opposite meanings: the people ironically
%		go from a desert (y[image]mn) to getting fat (y[image]mn).
%	}
}

\Verse{16}
{
	\hProto{יַקְנִאֻהוּ בְּזָרִים בְּתוֹעֵבֹת יַכְעִיסֻהוּ׃}
}
{
	\eDtrB{
		They made Him jealous with outsiders. With offensive
		things they made him angry.
	}
}
{
}

\Verse{17}
{
	\hProto{יִזְבְּחוּ לַשֵּׁדִים לֹא אֱלֹהַּ אֱלֹהִים לֹא יְדָעוּם חֲדָשִׁים מִקָּרֹב בָּאוּ לֹא שְׂעָרוּם אֲבֹתֵיכֶם׃}
}
{
	\eDtrB{
		They sacrificed to demons, a non-god, gods they hadn’t
		known; new ones, they came of late; your fathers hadn’t
		been acquainted with them.
	}
}
{
}

\Verse{18}
{
	\hProto{צוּר יְלָדְךָ תֶּשִׁי וַתִּשְׁכַּח אֵל מְחֹלְלֶךָ׃}
}
{
	\eDtrB{
		The rock that fathered you, you ignored, and you forgot
		{\God} who bore you.
	}
}
{
}

\Verse{19}
{
	\hProto{וַיַּרְא יְהֹוָה וַיִּנְאָץ מִכַּעַס בָּנָיו וּבְנֹתָיו׃}
}
{
	\eDtrB{
		And YHWH saw and rejected from His sons’ and His
		daughters’ angering.
	}
}
{
}

\Verse{20}
{
	\hProto{וַיֹּאמֶר אַסְתִּירָה פָנַי מֵהֶם אֶרְאֶה מָה אַחֲרִיתָם כִּי דוֹר תַּהְפֻּכֹת הֵמָּה בָּנִים לֹא אֵמֻן בָּם׃}
}
{
	\eDtrB{
		And He said, “Let me hide my face from them; I’ll see what
		their future will be. For they’re a generation of
		overthrows, children with no trust in them.
	}
}
{
%	\footnote{
%		Let me hide my face. See the comment on Deut 31:17. Footnote 32:20. I’ll
%		see what their future will be. Hebrew ’a[image][image]rît, often
%		understood as “end,” does not mean this in the sense of humankind’s
%		coming to an end. Rather, it refers to what will happen in the distant
%		future, the long run. There are two sides to the hiding of the face. It
%		is a fearful period of estrangement from their creator. It has been
%		called divine eclipse, Deus Absconditus, and “death” of God. But, after
%		living through it, humans are forced to grow up, to become more
%		responsible for their world. The words of the Song of Moses declare that
%		God is not simply hiding His face to bring His relationship with humans
%		to an end. There are rather two halves to the statement: “I’ll hide my
%		face from them; I’ll see what their future will be.” God gives humans
%		responsibility for their world. And the only way that they can be forced
%		to take that responsibility is if God is hidden. At whatever price (even
%		world wars, even the holocaust?!) and with whatever successes
%		(rebuilding Israel, conquering diseases, discovering secrets of the
%		creation of the universe), humans must grow up. And God will see what
%		their destiny will be. The Torah does not end with the natural
%		conclusions of the story: the promised land has not been reached, but it
%		is in sight. Everything lies in the future: finding a home, fulfillment
%		of the promises, bringing blessing to all the families of the earth. The
%		divine words resound: “I’ll see what their future will be.” The Torah
%		ends leaving us looking forward to what we can do and what we can be.
%	}
}

\Verse{21}
{
	\hProto{הֵם קִנְאוּנִי בְלֹא אֵל כִּעֲסוּנִי בְּהַבְלֵיהֶם וַאֲנִי אַקְנִיאֵם בְּלֹא עָם בְּגוֹי נָבָל אַכְעִיסֵם׃}
}
{
	\eDtrB{
		They made me jealous with no-god. They made me angry with
		their nothings. And I: I’ll make them jealous with
		no-people. I’ll make them angry with a foolish nation.
	}
}
{
}

\Verse{22}
{
	\hProto{כִּי אֵשׁ קָדְחָה בְאַפִּי וַתִּיקַד עַד שְׁאוֹל תַּחְתִּית וַתֹּאכַל אֶרֶץ וִיבֻלָהּ וַתְּלַהֵט מוֹסְדֵי הָרִים׃}
}
{
	\eDtrB{
		For fire has ignited in my anger and burned to Sheol at
		bottom and consumed land and its crop and set the
		mountains’ foundations ablaze.
	}
}
{
}

\Verse{23}
{
	\hProto{אַסְפֶּה עָלֵימוֹ רָעוֹת חִצַּי אֲכַלֶּה בָּם׃}
}
{
	\eDtrB{
		I’ll mass bad things over them. I’ll exhaust my arrows on
		them,
	}
}
{
}

\Verse{24}
{
	\hProto{מְזֵי רָעָב וּלְחֻמֵי רֶשֶׁף וְקֶטֶב מְרִירִי וְשֶׁן בְּהֵמֹת אֲשַׁלַּח בָּם עִם חֲמַת זֹחֲלֵי עָפָר׃}
}
{
	\eDtrB{
		sapped by hunger and devoured by flame. And bitter
		destruction and animals’ teeth I’ll let loose at them with
		venom of serpents of the dust.
	}
}
{
%	\footnote{
%		serpents of the dust. Or: things that crawl in the dust. This is
%		reminiscent of the story of Eden, in which the snake and its descendants
%		are condemned to be legless, going on their bellies and eating dust (Gen
%		3:14); and in which the descendants of the humans and the snake are
%		cursed with mutual enmity, so that the snake will bite the human’s heel
%		(3:15).
%	}
}

\Verse{25}
{
	\hProto{מִחוּץ תְּשַׁכֶּל חֶרֶב וּמֵחֲדָרִים אֵימָה גַּם בָּחוּר גַּם בְּתוּלָה יוֹנֵק עִם אִישׁ שֵׂיבָה׃}
}
{
	\eDtrB{
		Outside: a sword will bereave, and inside: terror, both
		young man and virgin, suckling with aged man.
	}
}
{
}

\Verse{26}
{
	\hProto{אָמַרְתִּי אַפְאֵיהֶם אַשְׁבִּיתָה מֵאֱנוֹשׁ זִכְרָם׃}
}
{
	\eDtrB{
		I’d say, ‘I’ll erase them. I’ll make their memory cease
		from mankind,’
	}
}
{
%	\footnote{
%		erase them … make their memory cease. To capture the full irony and
%		horror of these words, one must recall that this is what is supposed to
%		happen to Amalek, the people who symbolize Israel’s worst enemies: “wipe
%		out the memory of Amalek” (Deut 25:19).
%	}
}

\Verse{27}
{
	\hProto{לוּלֵי כַּעַס אוֹיֵב אָגוּר פֶּן יְנַכְּרוּ צָרֵימוֹ פֶּן יֹאמְרוּ יָדֵנוּ רָמָה וְלֹא יְהֹוָה פָּעַל כּל זֹאת׃}
}
{
	\eDtrB{
		if I didn’t fear the enemy’s anger, in case their foes
		would misread, in case they’d say, ‘Our hand was high, and
		it wasn’t YHWH who did all this!’
	}
}
{
}

\Verse{28}
{
	\hProto{כִּי גוֹי אֹבַד עֵצוֹת הֵמָּה וְאֵין בָּהֶם תְּבוּנָה׃}
}
{
	\eDtrB{
		For they’re a nation void of counsel, and there’s no
		understanding in them.
	}
}
{
}

\Verse{29}
{
	\hProto{לוּ חָכְמוּ יַשְׂכִּילוּ זֹאת יָבִינוּ לְאַחֲרִיתָם׃}
}
{
	\eDtrB{
		If they were wise they’d comprehend this; they’d grasp
		their future.
	}
}
{
}

\Verse{30}
{
	\hProto{אֵיכָה יִרְדֹּף אֶחָד אֶלֶף וּשְׁנַיִם יָנִיסוּ רְבָבָה אִם לֹא כִּי צוּרָם מְכָרָם וַיהֹוָה הִסְגִּירָם׃}
}
{
	\eDtrB{
		How could one pursue a thousand and two chase ten thousand
		if not that their rock had sold them, that YHWH had turned
		them over?”
	}
}
{
%	\footnote{
%		turned them over. This is the same word that is used to refer to turning
%		an escaped slave over to his master (Deut 23:16). The preceding colon,
%		which parallels this, refers to selling the people. That is, both lines
%		convey an image of God selling Israel back into slavery, when it was God
%		who had freed them from slavery in the first place. Like the last curse
%		of the list in Deuteronomy 28, it is the ultimate horror for Israel: to
%		return to slavery.
%	}
}

\Verse{31}
{
	\hProto{כִּי לֹא כְצוּרֵנוּ צוּרָם וְאֹיְבֵינוּ פְּלִילִים׃}
}
{
	\eDtrB{
		For their rock is not like our rock, though our enemies
		are the judges.
	}
}
{
}

\Verse{32}
{
	\hProto{כִּי מִגֶּפֶן סְדֹם גַּפְנָם וּמִשַּׁדְמֹת עֲמֹרָה עֲנָבֵמוֹ עִנְּבֵי רוֹשׁ אַשְׁכְּלֹת מְרֹרֹת לָמוֹ׃}
}
{
	\eDtrB{
		For their vine is from Sodom’s vine and from Gomorrah’s
		fields. Their grapes are poison grapes. They have bitter
		clusters.
	}
}
{
}

\Verse{33}
{
	\hProto{חֲמַת תַּנִּינִם יֵינָם וְרֹאשׁ פְּתָנִים אַכְזָר׃}
}
{
	\eDtrB{Their wine is serpents’ venom and cruel poison of cobras.}
}
{
}

\Verse{34}
{
	\hProto{הֲלֹא הוּא כָּמֻס עִמָּדִי חָתוּם בְּאוֹצְרֹתָי׃}
}
{
	\eDtrB{“Isn’t it stored with me, sealed in my treasuries?}
}
{
}

\Verse{35}
{
	\hProto{לִי נָקָם וְשִׁלֵּם לְעֵת תָּמוּט רַגְלָם כִּי קָרוֹב יוֹם אֵידָם וְחָשׁ עֲתִדֹת לָמוֹ׃}
}
{
	\eDtrB{
		Vengeance and recompense are mine, for the time their foot
		will slip; for the day of their ordeal is close and comes
		fast: things prepared for them.”
	}
}
{
}

\Verse{36}
{
	\hProto{כִּי יָדִין יְהֹוָה עַמּוֹ וְעַל עֲבָדָיו יִתְנֶחָם כִּי יִרְאֶה כִּי אָזְלַת יָד וְאֶפֶס עָצוּר וְעָזוּב׃}
}
{
	\eDtrB{
		For YHWH will judge His people and regret about His
		servants when He’ll see that their strength is gone and
		there’s none: held back or left alone.
	}
}
{
}

\Verse{37}
{
	\hProto{וְאָמַר אֵי אֱלֹהֵימוֹ צוּר חָסָיוּ בוֹ׃}
}
{
	\eDtrB{
		And He’ll say, “Where are their gods, the rock in whom
		they sought refuge,
	}
}
{
}

\Verse{38}
{
	\hProto{אֲשֶׁר חֵלֶב זְבָחֵימוֹ יֹאכֵלוּ יִשְׁתּוּ יֵין נְסִיכָם יָקוּמוּ וְיַעְזְרֻכֶם יְהִי עֲלֵיכֶם סִתְרָה׃}
}
{
	\eDtrB{
		who would eat the fat of their sacrifices, would drink the
		wine of their libations? Let them get up and help you! Let
		that be a shelter over you.
	}
}
{
}

\Verse{39}
{
	\hProto{רְאוּ עַתָּה כִּי אֲנִי אֲנִי הוּא וְאֵין אֱלֹהִים עִמָּדִי אֲנִי אָמִית וַאֲחַיֶּה מָחַצְתִּי וַאֲנִי אֶרְפָּא וְאֵין מִיָּדִי מַצִּיל׃}
}
{
	\eDtrB{
		See now that I, I am He, and there is no god with me. I
		cause death and give life. I’ve pierced, and I’ll heal.
		And there’s no deliverer from my hand,
	}
}
{
%	\footnote{
%		there is no god with me. It is possible that this might mean that other
%		gods exist but that they just do not happen to be present with YHWH at
%		this time, but that seems to me to be a stretch. Where would a god be
%		who is not in YHWH’s presence?! The plain sense of these words (and
%		other passages on which I have commented) is monotheistic. It comes in a
%		context of calling on gods to show themselves if they exist: “Where are
%		their gods.... Let them get up and help you!” And this poem shows
%		well-known linguistic signs of being an early composition. The
%		persistent view in biblical scholarship that monotheism is a late
%		development in biblical Israel (coming during or after the Babylonian
%		exile) is contrary to such explicit statements as this.
%	}
}

\Verse{40}
{
	\hProto{כִּי אֶשָּׂא אֶל שָׁמַיִם יָדִי וְאָמַרְתִּי חַי אָנֹכִי לְעֹלָם׃}
}
{
	\eDtrB{
		when I raise my hand to the skies, and I say, ‘As I live
		forever,
	}
}
{
}

\Verse{41}
{
	\hProto{אִם שַׁנּוֹתִי בְּרַק חַרְבִּי וְתֹאחֵז בְּמִשְׁפָּט יָדִי אָשִׁיב נָקָם לְצָרָי וְלִמְשַׂנְאַי אֲשַׁלֵּם׃}
}
{
	\eDtrB{
		if I whet the lightning of my sword, and my hand takes
		hold of judgment, I’ll give back vengeance to my foes and
		pay back those who hate me.
	}
}
{
}

\Verse{42}
{
	\hProto{אַשְׁכִּיר חִצַּי מִדָּם וְחַרְבִּי תֹּאכַל בָּשָׂר מִדַּם חָלָל וְשִׁבְיָה מֵרֹאשׁ פַּרְעוֹת אוֹיֵב׃}
}
{
	\eDtrB{
		I’ll make my arrows drunk with blood, and my sword will
		eat flesh from the blood of the slain and captured, from
		the head of loose hair of the enemy.’”
	}
}
{
}

\Verse{43}
{
	\hProto{הַרְנִינוּ גוֹיִם עַמּוֹ כִּי דַם עֲבָדָיו יִקּוֹם וְנָקָם יָשִׁיב לְצָרָיו וְכִפֶּר אַדְמָתוֹ עַמּוֹ׃}
}
{
	\eDtrB{
		Nations: cheer His people! For He’ll requite His servants’
		blood and give back vengeance to His foes and make
		atonement for His land, His people.
	}
}
{
}

\Verse{44}
{
	\hDtrB{וַיָּבֹא מֹשֶׁה וַיְדַבֵּר אֶת כּל דִּבְרֵי הַשִּׁירָה הַזֹּאת בְּאזְנֵי הָעָם הוּא וְהוֹשֵׁעַ בִּן נוּן׃}
}
{
	\eDtrB{
		And Moses came and spoke all the words of this song in the
		people’s ears, he and Hoshea son of Nun.
	}
}
{
%	\footnote{
%		Hoshea. Joshua.
%	}
}

\Verse{45}
{
	\hDtrA{וַיְכַל מֹשֶׁה לְדַבֵּר אֶת כּל הַדְּבָרִים הָאֵלֶּה אֶל כּל יִשְׂרָאֵל׃}
}
{
	\eDtrA{
		And Moses finished speaking all these things to all of
		Israel.
	}
}
{
}

\Verse{46}
{
	\hDtrA{וַיֹּאמֶר אֲלֵהֶם שִׂימוּ לְבַבְכֶם לְכל הַדְּבָרִים אֲשֶׁר אָנֹכִי מֵעִיד בָּכֶם הַיּוֹם אֲשֶׁר תְּצַוֻּם אֶת בְּנֵיכֶם לִשְׁמֹר לַעֲשׂוֹת אֶת כּל דִּבְרֵי הַתּוֹרָה הַזֹּאת׃}
}
{
	\eDtrA{
		And he said to them, “Pay attention to all the things that
		I testify regarding you today, that you’ll command them to
		your children, to observe and to do all the words of this
		instruction.
	}
}
{
}

\Verse{47}
{
	\hDtrA{כִּי לֹא דָבָר רֵק הוּא מִכֶּם כִּי הוּא חַיֵּיכֶם וּבַדָּבָר הַזֶּה תַּאֲרִיכוּ יָמִים עַל הָאֲדָמָה אֲשֶׁר אַתֶּם עֹבְרִים אֶת הַיַּרְדֵּן שָׁמָּה לְרִשְׁתָּהּ׃}
}
{
	\eDtrA{
		Because it’s not an empty thing for you, because it’s your
		life. And through this thing you’ll extend days on the
		land to which you’re crossing the Jordan to take
		possession of it.”
	}
}
{
%	\footnote{
%		it’s your life. Again we are reminded here at the end of the Torah that
%		these things are a path back to the tree of life that was lost at the
%		Torah’s beginning. See the comment on Deut 30:19.
%	}
}

\Verse{48}
{
	\hR{וַיְדַבֵּר יְהֹוָה אֶל מֹשֶׁה בְּעֶצֶם הַיּוֹם הַזֶּה לֵאמֹר׃}
}
{
	\eR{And YHWH spoke to Moses in this very day, saying,}
}
{
}

\Verse{49}
{
	\hR{עֲלֵה אֶל הַר הָעֲבָרִים הַזֶּה הַר נְבוֹ אֲשֶׁר בְּאֶרֶץ מוֹאָב אֲשֶׁר עַל פְּנֵי יְרֵחוֹ וּרְאֵה אֶת אֶרֶץ כְּנַעַן אֲשֶׁר אֲנִי נֹתֵן לִבְנֵי יִשְׂרָאֵל לַאֲחֻזָּה׃}
}
{
	\eR{
		“Go up to this mountain of Abarim, Mount Nebo, which is in
		the land of Moab, which is facing Jericho, and see the
		land of Canaan, which I’m giving to the children of Israel
		for a possession.
	}
}
{
}

\Verse{50}
{
	\hR{וּמֻת בָּהָר אֲשֶׁר אַתָּה עֹלֶה שָׁמָּה וְהֵאָסֵף אֶל עַמֶּיךָ כַּאֲשֶׁר מֵת אַהֲרֹן אָחִיךָ בְּהֹר הָהָר וַיֵּאָסֶף אֶל עַמָּיו׃}
}
{
	\eR{
		And die in the mountain to which you’re going up, and be
		gathered to your people, as Aaron, your brother, died in
		Mount Hor and was gathered to his people,
	}
}
{
}

\Verse{51}
{
	\hR{עַל אֲשֶׁר מְעַלְתֶּם בִּי בְּתוֹךְ בְּנֵי יִשְׂרָאֵל בְּמֵי מְרִיבַת קָדֵשׁ מִדְבַּר צִן עַל אֲשֶׁר לֹא קִדַּשְׁתֶּם אוֹתִי בְּתוֹךְ בְּנֵי יִשְׂרָאֵל׃}
}
{
	\eR{
		because you made a breach with me among the children of
		Israel at the waters of Meribah of Kadesh at the
		wilderness of Zin, because you didn’t make me holy among
		the children of Israel.
	}
}
{
}

\Verse{52}
{
	\hR{כִּי מִנֶּגֶד תִּרְאֶה אֶת הָאָרֶץ וְשָׁמָּה לֹא תָבוֹא אֶל הָאָרֶץ אֲשֶׁר אֲנִי נֹתֵן לִבְנֵי יִשְׂרָאֵל׃}
}
{
	\eR{
		Because you’ll see the land from opposite, but you shall
		not come there, to the land that I’m giving to the
		children of Israel.
	}
}
{
}
